% Slide 13: Hyper-Reduction: The Need for ECSW
\begin{frame}{Hyper-Reduction: The Need for ECSW}
    \begin{itemize}
        \item \textbf{Why Hyper-reduction?}: Galerkin ROM reduces DOFs ($15$), but online assembly scaling remains bound to the size of the full mesh ($N = 6,564$).
        \item Unreduced Galerkin ROM actually runs at \textbf{$0.92\times$ speedup} (slower than FOM) due to the projection assembly step.
        \item \textbf{Energy-Conserving Sampling \& Weighting (ECSW)}: Evaluates stiffness matrices over a sparse subset of active elements $E^* \subset E$, completely breaking the dependency on full mesh dimensions:
        \begin{equation*}
            \mathbf{K}_{r, \text{ECSW}}(\boldsymbol{\mu}) = \sum_{e \in E^*} w_e \boldsymbol{\Phi}_e^T \mathbf{K}_{e}(\boldsymbol{\mu}) \boldsymbol{\Phi}_e
        \end{equation*}
    \end{itemize}
\end{frame}

% Slide 14: ECSW Mathematical Formulation
\begin{frame}{ECSW Mathematical Formulation}
    \textbf{Non-Negative Least Squares (NNLS) Formulation:}
    \begin{equation*}
        \min_{\mathbf{w} \geq \mathbf{0}} \left\| \sum_{e=1}^{E} w_e \mathbf{z}_e - \mathbf{z}_0 \right\|_2^2
    \end{equation*}
    where:
    \begin{itemize}
        \item $\mathbf{z}_e = \text{vec}(\boldsymbol{\Phi}_e^T \mathbf{K}_e \boldsymbol{\Phi}_e)$ is the vectorized element stiffness contribution.
        \item $\mathbf{z}_0 = \text{vec}(\boldsymbol{\Phi}^T \mathbf{K}_{\text{full}} \boldsymbol{\Phi})$ is the target projected full-order stiffness matrix.
        \item The sparse solution restricts the active set $E^* = \{ e \mid w_e > 0 \}$ to a fraction of the total elements.
    \end{itemize}
\end{frame}

% Slide 15: ECSW Training and Coercivity Seeding
\begin{frame}{ECSW Training and Coercivity Seeding}
    \begin{columns}
        \begin{column}{0.6\textwidth}
            \textbf{Offline Training Strategy (Appendix B):}
            \begin{itemize}
                \item Training snapshots collected under $M=3$ loading configurations.
                \item \textbf{Coercivity Issue}: Standard NNLS selection may drop boundary elements, causing singular matrices.
                \item \textbf{Pre-seeding Solution}: Element indices on boundary Dirichlet constraints and high stress-gradient notch zones are pre-seeded into the active set $E^*$.
            \end{itemize}
        \end{column}
        \begin{column}{0.38\textwidth}
            \centering
            \begin{tikzpicture}[scale=0.8]
                % Boundary
                \draw[fill=red!70!black] (-2.2, -0.4) rectangle (-2.0, 0.4);
                \node[left, font=\tiny] at (-2.2,0) {Dirichlet Seed};
                
                % Beam
                \draw[draw=navyblue, fill=cyansec!10, thick] (-2.0, -0.4) -- (2.0, -0.4) -- (2.0, 0.4) -- (-2.0, 0.4) -- cycle;
                
                % Active bands
                \fill[blue!40!purple, opacity=0.7] (-1.8, -0.4) rectangle (-1.6, 0.4);
                \fill[blue!40!purple, opacity=0.7] (-0.2, -0.4) rectangle (0.0, 0.4);
                \fill[blue!40!purple, opacity=0.7] (1.4, -0.4) rectangle (1.6, 0.4);
                
                \node[font=\tiny] at (0,0) {Active Elements $E^*$};
            \end{tikzpicture}
        \end{column}
    \end{columns}
\end{frame}

% Slide 16: Dynamic Weight Calibration
\begin{frame}{Dynamic Weight Calibration}
    \begin{itemize}
        \item \textbf{Problem}: Pullback to $\hat{\Omega}$ distorts B-spline weights under extreme notch depths ($r > 0.4$).
        \item \textbf{Solution}: Real-time weight calibration factor $\gamma$ applied to the active set:
        \begin{equation*}
            w_e^*(\boldsymbol{\mu}) = \gamma(\boldsymbol{\mu}) w_e
        \end{equation*}
        \item \textbf{Calibration Parameter}:
        \begin{equation*}
            \gamma(\boldsymbol{\mu}) = 1.0 + c_1 r^2 + c_2 \left(\frac{\mu - \mu_0}{\mu_0}\right)^2
        \end{equation*}
        \item Corrects projection distortions, maintaining $L_2$ error below $0.06\%$ across the entire parameter space.
    \end{itemize}
\end{frame}

% Slide 17: Quantitative Results: Solver Speedup
\begin{frame}{Quantitative Results: Solver Speedup}
    \begin{table}
        \centering
        \caption{Performance metrics benchmarked on the notched cantilever geometry ($E = 3,200$, $N = 6,564$ DoFs).}
        \label{tab:performance_results_extended}
        \begin{tabular}{lcccc}
            \toprule
            \textbf{Solver Strategy} & \textbf{Active DoFs / Cells} & \textbf{Relative $L_2$ Error (\%)} & \textbf{Assembly Time} & \textbf{Speedup} \\
            \midrule
            Full Order (FOM)        & $6,564$ DoFs / $3,200$ Cells & Baseline & $142.1$ ms & $1.0\times$ \\
            Galerkin ROM (Unreduced)& $15$ DoFs / $3,200$ Cells    & $0.012\%$   & $154.5$ ms & $0.92\times$ \\
            \textbf{ECSW ROM (Hyper-reduced)}& \textbf{15 DoFs / 28 Cells}  & \textbf{0.054\%} & \textbf{2.2 ms} & \textbf{64.5$\times$} \\
            \bottomrule
        \end{tabular}
    \end{table}
    \begin{itemize}
        \item ECSW ROM evaluates only \textbf{$28$ active cells} instead of $3,200$, reducing assembly time to \textbf{$2.2\text{ ms}$}.
        \item Delivers a speedup of \textbf{$64.5\times$} while maintaining sub-millimeter precision.
    \end{itemize}
\end{frame}

% Slide 18: Quantitative Results: Spectral/Modal Conservation
\begin{frame}{Quantitative Results: Spectral/Modal Conservation}
    \textbf{Spectrum Conservation Check:}
    \begin{itemize}
        \item \textbf{Fundamental Frequency}: $\omega_{n,\text{FOM}} = 1.108\text{ rad/s}$ vs. $\omega_{n,\text{ECSW}} = 1.107\text{ rad/s}$ ($0.09\%$ deviation).
        \item Conservations of energy-based frequencies ensures structural stability during dynamic transients.
    \end{itemize}
    \vspace{0.2cm}
    \centering
    \begin{tikzpicture}[scale=0.7]
        \begin{axis}[
            width=0.85\textwidth,
            height=0.25\textwidth,
            xlabel={Time $t$ [s]},
            ylabel={$U_y$ [mm]},
            xmin=1.2, xmax=5.2,
            ymin=-0.5, ymax=0.5,
            grid=both,
            grid style={line width=.1pt, draw=gray!10},
            major grid style={line width=.2pt, draw=gray!20},
            axis lines=left,
            x tick label style={font=\tiny},
            y tick label style={font=\tiny},
        ]
            \addplot[color=navyblue, thick] coordinates {
                (1.2, 0.4014) (1.4, 0.3910) (1.6, 0.4077) (1.8, 0.3994) (2.0, 0.3256) 
                (2.2, 0.2037) (2.4, 0.0847) (2.6, -0.0040) (2.8, -0.0794) (3.0, -0.1703) 
                (3.2, -0.2767) (3.4, -0.3688) (3.6, -0.4197) (3.8, -0.4306) (4.0, -0.4224) 
                (4.2, -0.4072) (4.4, -0.3764) (4.6, -0.3152) (4.8, -0.2231) (5.0, -0.1163) 
                (5.15, -0.0375)
            };
            \addplot[color=cyansec, dashed, thick] coordinates {
                (1.2, 0.4018) (1.4, 0.3915) (1.6, 0.4072) (1.8, 0.3990) (2.0, 0.3260) 
                (2.2, 0.2040) (2.4, 0.0849) (2.6, -0.0038) (2.8, -0.0790) (3.0, -0.1700) 
                (3.2, -0.2760) (3.4, -0.3680) (3.6, -0.4190) (3.8, -0.4300) (4.0, -0.4220) 
                (4.2, -0.4070) (4.4, -0.3760) (4.6, -0.3150) (4.8, -0.2230) (5.0, -0.1160) 
                (5.15, -0.0374)
            };
        \end{axis}
    \end{tikzpicture}
\end{frame}
